\solution{Магнитные цепи и трансформатор}{10.0}

\tasksol Это --- стандартная задача о магнитном поле бесконечно длинного соленоида. По теореме Ампера для прямоугольного контура длины $l$ мы имеем

\begin{eq}[points=0.2]
    H\cdot l = NI,
\end{eq}
и с учётом $B=\mu\mu_0 H$ внутри сердечника получаем

\begin{eq}[points=0.1, numpoints=0.1]
B = \frac{\mu\mu_0 NI}{l}=\qty{1}{\T}.
\end{eq}

\criterion[type=remark]{Если участник сразу записал \eqref{1} через $B=\mu\mu_0 H$, даётся балл за \eqref{1} тоже}

\tasksol Магнитный поток дан как

\begin{eq}[points=0.2]
    \Phi = B\cdot \pi r^2.
\end{eq}

Объединяя с \eqref{2}, получаем

\begin{eq}[points=0.2, numpoints=0.2]
    R_m = \frac l{\mu\mu_0\cdot\pi r^2}=\qty{6.33e5}{\A\per\Wb}.
\end{eq}

Здесь эквивалентны единицы $\unit{\A\per\Wb}=\!\unit{\A\per\T\per\m\squared}=\unit{\per\H}$.

\tasksol Индуктивность равна потокосцеплению $\Psi$ на единицу пускаемого тока. Здесь

\begin{eq}[points=0.2]
    \Psi=N\Phi,
\end{eq}

так что

\begin{eq}[points=0.2, numpoints=0.2]
    L = \frac{N^2}{R_m} = \frac{\mu\mu_0\pi r^2N^2}l =\qty{63.2}{\mH}.
\end{eq}

\tasksol Для рассмотрения граничных условий нужно воспользоваться двумя фактами: соленоидальность магнитного поля $\nabla\cdot\vec B=0$, которая сводится к тому, что при применении теоремы Гаусса на небольшой цилиндр, смотрящий обоими торцами в разные среды, мы получаем \textit{сохранение нормальной компоненты вектора} $\vec B$; а также \textit{сохранение тангенциальной компоненты вектора} $\vec H=\vec B/\mu\mu_0$ при использовании $\oint\vec H\cdot d\vec l=0$ для прямоугольного контура, пересекающего перпендикулярно границам раздела. Сохранение нормальной компоненты вектора $\vec B$ даёт

\begin{eq}[points=0.3]
    B_0\cos\alpha=B\cos\beta,
\end{eq}

а сохранение тангенциальной компоненты вектора $\vec H$ даёт (поскольку $\mu=1$ для воздуха)

\begin{eq}[points=0.3]
    B_0\sin\alpha = \frac B\mu \sin\beta.
\end{eq}

Решая вместе \eqref{7}--\eqref{8}, получаем

\begin{eq}[points=0.2]
    \tan\beta = \mu \tan\alpha,
\end{eq}

т.е.

\begin{eq}[points=0.1, prefix={Численное значение}]
    \beta \approx \ang{89.83},
\end{eq}

а также

\begin{eq}[points=0.2, numpoints=0.1]
    \frac B{B_0}=\sqrt{\cos^2\alpha+\mu^2\sin^2\alpha} \approx 347.
\end{eq}

Видно, что почти для любых углов $\alpha$, кроме $\alpha\ll1$ в радианной мере, мы всегда имеем $\beta$ очень близкий к 90 градусам, а $B_0$ заметно меньше $B$, что можно безопасно пренебрегать утечкой потока в подобных задачах.

\tasksol Используем теорему Ампера вдоль центра сердечника:

\begin{eq}[points=0.2]
    H_1\cdot(4a-d) + H_2\cdot d = NI,
\end{eq}

где $H_1$ --- напряжённость магнитного поля внутри сердечника, а $H_2$ --- в воздушном зазоре. Иначе говоря, ввиду сохранения $B$,

\begin{eq}[points=0.3]
    B = \mu_0 H_2 = \frac{\mu_0}\mu H_1.
\end{eq}

Итак, мы получаем для потока

\begin{eq}[points=0.3]
    \Phi = \frac{\mu_0 NI S}{d + \mu(4a-d)}.
\end{eq}

Однако если сравнить с формулой \eqref{4}, мы видим, что можно использовать

\begin{eq}[points=0.1]
    R_{m,1} = \frac{4a-d}{\mu\mu_0S},
\end{eq}

\begin{eq}[points=0.1]
    R_{m,2} = \frac d{\mu_0S},
\end{eq}

так что это аналогично записи $R_m=R_{m,1}+R_{m,2}$ и применения формулы (1) из условия задачи.

\tasksol Задав $R_{m,1}=R_{m,2}$, получим

\begin{eq}[points=0.2, numpoints=0.2]
    d = \frac{4a}{\mu+1}\approx\frac{4a}{\mu} = \qty{0.4}{\mm}.
\end{eq}

Эта часть наглядно демонстрирует, что магнитное сопротивление преимущественно концентрируется в воздушных зазорах!

\tasksol%
\criterion{В пунктах \solref{7}, \solref{10}--\solref{11}, связанные с законами магнитных цепей критерии (\eqref{18}, \eqref{20}, \eqref{27}--\eqref{29}, \eqref{33}--\eqref{34}) принимаются при наличии аналогичных выражений}%
%
Сперва определим коэффициент взаимной индукции, которая является лишь геометрическим свойством двух катушек. Поэтому для простоты положим, что есть лишь ток $I_1$, создаваемый $V(t)$, а $I_2=0$. Поток, создаваемый первичной обмоткой, равен

\begin{eq}[points=0.2]
    \Phi(t) = \frac{N_1 I_1(t)}{R_m},
\end{eq}

где $R_m$ определён из \eqref{15} для $d=0$. Это создаёт потокосцепление на вторичную обмотку

\begin{eq}[points=0.2]
    \Psi_2(t) = N_2\Phi(t),
\end{eq}

так что из определения коэффициента взаимной индукции $M = \Psi_2/I_1$ получаем

\begin{eq}[points=0.2]
    M = \frac{N_1N_2}{R_m} = \frac{\mu\mu_0N_1N_2S}{4a}.
\end{eq}

Теперь посмотрим на реальную цепь и определим эквивалентное сопротивление. По закону Фарадея

\begin{eq}[points=0.2]
    V(t) = -N_1\frac{d\Phi}{dt},
\end{eq}

а напряжение $V_2=I_2R$ создаётся через тот же закон $V_2=-N_2\dot\Phi$. Мы видим, что следовательно

\begin{eq}[points=0.2]
    I_2 R = \frac{N_1}{N_2}V(t).
\end{eq}

Поток определяется законом Ома для магнитной цепи: мы видим, что первичная и вторичная обмотки делают вклад в магнитодвижущую силу (МДС), однако сопротивление цепи, как было указано в подсказке, слишком мало, так что

\begin{eq}[points=0.2]
    N_1I_1+N_2I_2 = 0.
\end{eq}

Отсюда можем выразить силу тока $I_1$. Решая \eqref{21}--\eqref{23}, получаем

\begin{eq}[points=0.3, override={R'=R\ab(\frac{N_1}{N_2})^2}]
    I_1 =-\frac{V(t)}{R}\ab(\frac{N_2}{N_1})^2\implies R'=R\ab(\frac{N_1}{N_2})^2.
\end{eq}

Знак минус показывает, что сила тока сдвинута по фазе на $\pi$ относительно переменного напряжения $V(t)$.

\tasksol Как видно из определения о магнитном сопротивлении, магнитный поток есть отношение МДС к магнитному сопротивлению. Напряжённость электрического поля есть результат действия электродвижущей силы по закону Фарадея $\oint\vec E\cdot d\vec l =\mathcal E$, что аналогично теореме Ампера $\oint\vec H\cdot d\vec l = NI$. Магнитная проницаемость связывает напряжённость поля $\vec H$ с плотностью потока магнитной индукции $\vec B$, то есть последнее имеет аналогию с плотностью тока $\vec j$, связанной с электрическим полем соотношением $\vec j =\sigma \vec E$. Итого мы получаем

\begin{center}
    \begin{tabular}{|c|c|}\hline
         \textbf{Магнитная цепь}& \textbf{Электрическая цепь} \\\hline
         Магнитный поток $\Phi$ &Сила тока $I$\\\hline
         Напряжённость магнитного поля $\vec H$&Напряжённость электрического поля $\vec E$ \\\hline
         Абсолютная магнитная проницаемость $\mu\mu_0$&Удельная проводимость $\sigma\equiv1/\rho$\\\hline
    \end{tabular}
\end{center}

\criterion[points=0.1]{Записано: <<Сила тока $I$>>}
\criterion[points=0.2]{Записано: <<Напряжённость магнитного поля $\vec H$>>}
\criterion[points=0.2]{Записано: <<Удельная проводимость $\sigma$ или $1/\rho$>>}

\textit{Примечание жюри: в приоритет ставится написанное участником словесное название величины, и лишь затем его буквенная запись. При отсутствии записанного термина ставится балл за аналогию только если буквенная запись в точности совпадает с авторским.} Стоит отметить, что аналогии идут \textit{исключительно} в контексте электрических/магнитных цепей. Поэтому, к примеру, нельзя говорить о том, что диэлектрическая проницаемость $\varepsilon$ аналогична магнитной проницаемости $\mu$, хоть в другом контексте эта аналогия тоже верна.

\tasksol Мы имеем формулу плотности магнитной энергии

\begin{eq}[points=0.2]
    w_m = \frac{BH}{2} = \frac{B^2}{2\mu\mu_0}=\frac{\mu\mu_0 H^2}{2}.
\end{eq}

Если рассмотрим участок длины $l$ и площади сечения $S$, то $W_m=w_m Sl$, и принимая во внимание \eqref{4} и аналогию в \solref{8}, заключаем

\begin{eq}[points=0.3, override={W_M = \frac{\mathcal F^2}{2R_m}}]
    W_M = \frac{\mathcal F^2}{2R_m}=\frac12\Phi^2R_m.
\end{eq}

\pic{12}R{0pt}{0.45}{figures/3.7.png}

\tasksol Можно нарисовать эквивалентную электрическую цепь и увидеть, что МДС фазы и нуля совместно создают поток по часовой стрелке на внешнем контуре. Здесь $R_m=\rho_m\cdot l$ --- сопротивление участка длины $l$ (вся фигура имеют суммарную длину $7l$). При $I_1=I_2$ мы бы имели нулевой поток на центральном ``мостике'', а при ненулевом значении есть разность потоков $\Phi_1-\Phi_2$, которые протекают через данную перемычку. Вторичная обмотка лишь ``реагирует'' на приходящий поток, поэтому в эквивалентной цепи этот участок (вторичная обмотка, магнитный пускатель $M$, и сопротивление $R$) остаются как есть и технически не являются элементами магнитной цепи. Однако теперь мы можем записать второе правило Кирхгофа для, к примеру, внешнего контура

\begin{eq}[points=0.3]
    NI_1(t) + NI_2(t) = 3R_m(\Phi_1+\Phi_2)
\end{eq}

и для, к примеру, левого контура

\begin{eq}[points=0.3]
    NI_1(t)+N_0I_M(t) = 3R_m\cdot\Phi_1 + R_m\cdot (\Phi_1-\Phi_2).
\end{eq}

Решением \eqref{27}--\eqref{28} получаем для $\Delta\Phi=\Phi_1-\Phi_2$

\begin{eq}[points=0.3]
    \Delta\Phi = \frac{N(I_1(t)-I_2(t))}{5R_m}.
\end{eq}

Поскольку сказано, что магнитный пускатель не имеет импеданса, всё напряжение на вторичной обмотке уходит на резистор, задавая ток $I_M$. Через закон Фарадея получим

\begin{eq}[points=0.1]
    I_M(t) = -\frac{N_0}R\frac{d}{dt}\Delta\Phi,
\end{eq}

(знак минус необязателен, так как нам интересна амплитуда), и применяя на \eqref{29}, с учётом того, что $\frac{d}{dt}I(t) = 2\pi f\cdot I(t+\pi/2)$, получаем для амплитудного значения

\begin{eq}[points=0.3, numpoints=0.2]
    I_M = \frac{2\pi fNN_0(I_1-I_2)}{\sqrt{(5R\rho_m l)^2+(4\pi fN_0^2)^2}}=\qty{11.8}{\mA}.
\end{eq}

\tasksol Пусть $x$ --- сжатие пружин при включении $\mathcal F$. В таком случае в зазоре останется объём

\begin{eq}[points=0.1]
    V=\pi a^2(d-x).
\end{eq}

Так как $\mu\gg1$, из \solref{9} мы видим, что магнитная энергия пропорциональна $R_m$ для фиксированного потока, а из \solref{6} магнитное сопротивление создаётся преимущественно воздушным зазором, можно с высокой точностью положить, что большинство магнитной энергии запасается именно в воздушном зазоре. Для пространства между сердечниками \textbf{(b)} и \textbf{(c)} мы имеем

\begin{eq}[points=0.1]
    W_m = \frac{\mathcal F^2}{2R_m},
\end{eq}

где $R_m$ равен

\begin{eq}[points=0.4]
    R_m = \frac{d-x}{\mu_0 \pi a^2}.
\end{eq}

\criterion[type=remark]{Участник может записать вклады в энергию от сердечников или боковых зазоров; балл ставится если в их выражении имеется правильное слагаемое \eqref{33}--\eqref{34}}

Виртуальные перемещения $x$ создают градиент энергии, что проявляется в пондеромоторной силе

\begin{eq}[points=0.2]
    F=-\frac{dW_m}{dx},
\end{eq}

или же из \eqref{33},
$$
F(x) = \frac{\mu_0\pi a^2\mathcal F^2}{2(d-x)^2},
$$
то есть $F$ действует в сторону уменьшения $x$. Вообще говоря, магнитное поле существует не только в воздушном зазоре, но и в других частях пространства, однако геометрия этих областей не изменяется при перемещении якоря. Следовательно, соответствующая энергия не зависит от $x$ и не влияет на силу $F$. Итак, условием замыкания является

\begin{eq}[points=0.2]
    F(x_c) \geq kx_c.
\end{eq}

Заметим, что приравнивание изменения магнитной энергии с энергией пружины $kx^2/2=\Delta W_m$ здесь неприменимо. Такое соотношение описывает максимальное отклонение якоря при свободном движении под действием магнитной силы. В нашей задаче требуется определить условие срабатывания, т.е. момент, когда магнитная сила становится не меньше силы пружины. Поэтому необходимо сравнивать силы $F(x)$ и $kx$, а не их энергии. Итак, дифференцируя \eqref{33} согласно \eqref{35} и подставляя в \eqref{36}, получаем

\begin{eq}[points=0.5, override={\mathcal F_0 = \frac{d-x_c}a\sqrt\frac{2kx_c}{\mu_0\pi}}]
    \mathcal F_0 = \frac{d-x_c}a\sqrt\frac{2kx_c}{\mu_0\pi}=\qty{15.75}{\A}.
\end{eq}

Применяя $\mathcal F_0=NI_M$, получим численно

\begin{eq}[points=0.2, prefix={Численное значение}]
    I_M = \qty{11.25}{\mA}.
\end{eq}

Видно, что \eqref{31} незначительно превосходит \eqref{38} (ток, создаваемый дифференциальным трансформатором, немного превышает порог срабатывания), так что при данных условиях должен сработать предохранитель.

\criterion[points=0.1]{Показано, что $I_M<\qty{11.8}{\mA}$}
